\documentclass[leqno, a4paper, 12pt]{jsarticle}
\usepackage{amssymb}
\usepackage{amsthm}
\usepackage{amsmath}
\usepackage{comment}
\usepackage{url}

\usepackage{enumitem}
%\usepackage{showkeys}


%定理環境 thmはあまり使わずpropをなるべく使う。
%主張は基本的に証明の中で使い、そのため番号を付けない。
%注意環境を付けてみた。
\newtheorem{thm}{定理}
\newtheorem{prop}[thm]{命題}
\newtheorem{cor}[thm]{系}
\newtheorem{lem}[thm]{補題}
\newtheorem{df}[thm]{定義}
\newtheorem{exam}[thm]{例}
\newtheorem{fact}[thm]{事実}
\newtheorem*{claim}{主張}
\newtheorem*{cau}{注意}
\newtheorem*{rmk}{注意}
\renewcommand{\proofname}{証明}
%矢印たち。
%\toは\rightarrowと同じ。\getsは\leftarrowと同じ。同値は\iff
\newcommand{\To}{\Longrightarrow}
\newcommand{\Gets}{\LongLeftarrow}
%黒板太字
\newcommand{\nn}{\mathbb{N}}
\newcommand{\zz}{\mathbb{Z}}
\newcommand{\qq}{\mathbb{Q}}
\newcommand{\rr}{\mathbb{R}}
\newcommand{\cc}{\mathbb{C}}
%無限大
\newcommand{\mgn}{\infty}
%差集合は\setminusで統一する。等号の否定は\neq
%真部分集合は\subsetneqq　偏微分は\partial
%ディスプレイスタイル。\dfracでディスプレイ分数
\newcommand{\dpst}{\displaystyle}

\title{論文メモ
\large{\\--- ノンジャンル ---}
}
\author{ナンブ　キトラ}
\date{\today}
\begin{document}
\maketitle

本棚を整理したいので，
溜まっている論文のメモを書いて未来への手紙とすることにする．



なんかジャンル分けできなかった論文を集めた束．


\begin{enumerate}[label=(\arabic*)]

\item 
論文
\cite{MR0277667}
は局所コンパクトな従順群の特徴づけをしていると思う，
左不変平均を使って．


\item 
論文
\cite{MR1501735}
は有名な
ホイットニーの拡張定理の論文だが，
おそらく私の手元にあるのはそれの
アナウンスメントである．


\item 
論文
\cite{MR0697498}
は
可微分多様体の間の
可微分写像の空間に
ホイットニーの強位相を入れた空間は
ベール空間であることを証明している．



\item 
論文
\cite{Sierpinski1917}
はM. Borelによる
正規数の話の初等的証明みたいなタイトルだが，
フランス語がよくわからないのでわかんなかった．


\item 
論文
\cite{MR2757537}
はウルトラプロダクトの解説である．



\item 
論文
\cite{freiwald1973cardinalities}
では距離関数の完備化が扱われている．
距離空間$X$の濃度を$\kappa$
とすると完備化の濃度は$\kappa$
かもしくは$\kappa^{\aleph_{0}}$
らしいが,
そのそれぞれの場合において
$X$の位相に何が起こるのかを調べている．



\item 
論文
\cite{MR0023881}
は基数の冪の話である．
ちょっとよくわかんなかった．
基数の冪の話でAnn. of Math.
に載ってるのが珍しくてプリントしたのかも知れない．



\item 

論文
\cite{MR1506850}
はユークリッド幾何を
当時風に書いたようなもののような気がする．
確かユークリッド空間への埋め込み可能性が書かれていたようなそうでもないような．


\item 
論文
\cite{MR1249357}
はサーストンが書いた数学における証明と
進歩に関するエッセイだと思う．
多分なんか数学の発展とかのいいお話が書いてあんだと思う．



\item 
論文
\cite{MR0195099}
では
可微分多様体の可微分性を上げる話をしていると思う．
ホイットニーもそういうことやっていると思うが，
境界つき多様体でもやってるようなので
こっちの方がいいかも知れない．



\item 

論文
\cite{MR1511870}
のタイトルはnowhere dense な完全集合に関する注意だと書いてあると思うが内容はよくわからず．
多分なんかカントール集合の話をしていて何か
反例を作っているような気がするがよくわからず．



\item 

論文
\cite{MR3941145}
では
距離空間$K$であって
$\dim K^{2}=2\dim K$
が成り立つものをbasic type, 
$\dim K^{2}=2\dim K-1$
が成り立つものを
exceptional type 
とよんで
これらの$K$から誘導される
$K^{\zz}$の平均次元がそれぞれ
$\dim K$
と
$\dim K-1$
になることを証明しており大変面白い論文だと思う．
平均次元というのはグロモフによって導入された
位相力学系的な次元でそういうものと
ジェネトポ的なものが交わり面白いと思う．



\item 
論文
\cite{MR0476524}
はバリバリの記述集合論だと思う．
なんかボレル集合から生成される同値関係について
その同値類の個数が
ほんのちょっと連続体仮説っぽくなってることを証明している．



\item 
論文
\cite{MR0565886}
はユークリッド空間の領域上で定義された関数が
単射になる十分条件を賢く探索している．
なんだかいい話だと思うし，解析と距離関数みたいな議論が
交わっているような気がするのでよく読んだ方がいい気がするが時間がない．




\item 
論文
\cite{MR2904087}
はindependent family とresolvability
の話である．
この二つ関係あるのかと思って印刷したが
ゴリゴリの集合論だったのでよく読めず．


\item 

1953年の
論文
\cite{MR0063022}
では距離化可能空間に第二可算性などの仮定を置いて距離関数の拡張問題について研究して証明しているが，
残念ながらこのような結果はより一般的に
1930年にハウスドルフによって
証明されている．



\item 
論文
\cite{MR1171783}
はRadoの選択原理を
色分けなどに応用する話である．



\item 
プレプリン
\cite{noble2021yet}
はなんだろうこれ．
NobleさんというのはNobleの定理(任意の基数$\mathfrak{m}$について$X^{\mathfrak{m}}$が正規ならば$X$はコンパクトハウスドルフ)の
Nobleさんだと思うが書き出しが1971年にアナウンスした
チコノフの定理の良い証明をこのプレプリントで
紹介しているらしく，アナウンスから50年かかって定理
のベールが剥がれたということらしい．
Nobleらしく(?)，射影を使ってチコノフの定理を証明している．



\item 
論文
\cite{torrome2020quotient}
はよくわかんなかった．
$\zz$
の商群
に距離を入れてそれでグロモフ・ハウスドルフ距離を計算しているようだがよくわからず．


\item 
Whitneyの拡張定理といえば，
ユークリッド空間の閉集合上の「可微分関数」を
全体空間の可微分関数に拡張する定理を思い浮かべるが，
1951年の論文\cite{nakaoka1951whitney}
によると，
$Y$を単連結な弧状連結空間で$\pi_{2}(Y)$が有限性性なものとすると，任意の四次元複体$K$と任意の$f\colon K^{2}\to Y$は$K$全体に拡張されるための必要十分条件をWhitneyが与えたらしい．ここで$K^{2}$は多分$K$の二次元スケルトンだと思う．多分特性類とかの話の萌芽だと思う．


\item 
論文
\cite{wa1}
では
カントールによらない実数の非可算性の証明
の話が載っている．
また対角線論法や区間縮小法の歴史についても述べられており，
読んでいてとても楽しいと思う．
付録では色々な実数の非可算性の証明が紹介されており
これも楽しい．



\item 
論文
\cite{heath1973monotonically}
はmonotonically normal space
(単調正規空間)
の論文．
なんかGoogle Scholar で調べると
ドマイナーな話題なのに
260回も引用されているらしい．
この単調正規性は
一般化された距離空間の話題とも関係あるようだ．
とりあえず単調正規性について知りたければこの論文を読めば良いと思う．


論文
\cite{steen1970direct}
はLOTSが
遺伝的族正規であることの直接証明らしい．
論文
\cite{cater2006simple}
はこの事実の簡単な証明らしい．



ちなみに\cite{steen1970direct}は1970年で
\cite{heath1973monotonically}は1973年である．
論文\cite{heath1973monotonically}以後はLOTSが
遺伝的族正規を示す際には以下の論法が可能になった．
つまりLOTSは単調正規であり，
一般に
単調正規性から遺伝的族族正規も従うので
LOTSも遺伝的族正規である．
ちなみにLOTSが族正規であることの証明は
なぜか電波通信の可算パラコンパクトとDowker空間の記事で
紹介している(https://concious4410.hatenablog.com/entry/2019/11/29/182755)．
単調正規性から可算パラコンパクトが従うからついでに紹介しているのだと思う．

単調正規空間の論文は他にもある
→
\cite{MR3414874}, 
\cite{MR1193194}, 
\cite{MR1425937},
\cite{MR2356193}．
また論文
\cite{MR2575894}
はGO 空間が単調正規であることの
直接証明らしい．



\item 
論文
\cite{conopeconfo}
は等角写像の話だが，
ちょっとよくわかんなかった．


\item 
ホワイトヘッドの
1935年の
論文
\cite{MR1503245}
は可微分多様体に
いわゆるgood cover
が存在するという話の
原論文のような気がする．


\item 
次の論文たちは
ホイットニーの拡張定理の延長上の話だと思う
→
\cite{MR2233851}, 
\cite{MR2339466}, 
\cite{MR2134225}, 
\cite{10.1093/qmath/os-7.1.1}

またホイットニー自身の論文も参照のこと
→
\cite{MR0009785},
\cite{MR1501875},
\cite{MR1562803},
\cite{MR1503173},
\cite{MR1501749}. 

論文
\cite{MR1191578}
は一変数のホイットニーの拡張定理の
構成的証明らしい．


\item 

次の論文たちは
連続写像が常に一様連続になる空間に関するものであると思う
→
\cite{MR0184200}, 
\cite{MR0619571}, 
\cite{MR0801222}, 
\cite{MR0837486}, 
\cite{MR0810180},
\cite{MR0099023}, 
\cite{MR0966244}, 
\cite{MR0106448}, 
\cite{MR0672468}. 



\item 
Kuratowski--Mr\'{o}wka
の定理に関する論文は次のとおりである
→
\cite{MR0250275}, 
\cite{MR0203701},
\cite{MR2584761}, 
\cite{MR0328865}．
また
論文
\cite{MR0117704}
と
論文
\cite{Kuratowski1931}
は原論文である．

\item 

Hewittの論文
\cite{MR0026239}
は実コンパクトを定義した論文のような気がする．
関数環考えて素イデアル考えるとかそういうやつ．


\item 
論文
\cite{MR1717184}
は位相空間のタイリングの話らしいがよくわからなかった．



\item
論文
\cite{wa2}
は凸解析における不動点定理の話である．




\item 
論文
\cite{wa4}
は初等部分モデルを使った位相空間論の話である．
論文\cite{wa3}は位相次元の話である．


\item
論文
\cite{MR0111816}
はダニエル積分の話である．

\item 
論文
\cite{MR0062730}
では0次元コンパクト環
が多分profiniteみたいな
話をしている．

\item 
論文
\cite{MR0069247}
では$\lambda$が極限順序数であることと
任意の$2\le \xi<\omega$
を満たす順序数
$\xi$
について
$\lambda=\xi\lambda$
が成り立つことが同値であることを証明している．


\item 
族正規というのはBingが導入した概念だが，
論文
\cite{MR0063018}
では
正規性と可算族正規が同値であることを証明している．


\item 
論文
\cite{MR1556955}
はバーコフの，
位相群が第一可算だったら距離化可能っていうやつの原論文だと思う．


\item 

論文
\cite{MR0029152}
は無限直積空間の弱位相の話らしいがよくわからず．

\item 
論文
\cite{MR0242120}
と
\cite{MR0233333}
は距離空間のパラコンパクト性の証明だが，
今ではM. E. Rudinの証明が
有名である．ちょっと悲しい．

\item 
論文
\cite{Montgomrey1935}
は非可分な距離空間でのボレル階層の話だと思う．
この論文は何か大事な論文だったような気がするが，
どのように大事であったのか思い出せない．
論文
\cite{MR1367597}
もそういう非可分なボレル階層というか
記述集合論みたいな話．

\item 
論文
\cite{MR1502556}
はなんだろうよくわかんない．


\item 
論文\cite{MR0808727}
は写像があったときにそれをいい感じに径数付け
(parametrization)できるみたいな話だが，
多分 para の部分と
metrizationの部分だけ見て印刷したような気がする．
あとE. Michaelが著者の中に入っているので
パラコンパクトとか距離化可能定理っぽい話だとかなり間違えたのだと思う．



\item 
論文
\cite{Mathur1971}
はコンパクト性を弱めた変種の論文である．
可算コンパクトとか．特に184ページの同値性の図はわけわからんことになってて面白い．

論文
\cite{MR0314002}
は可算コンパクト性とかの話．


\item 
論文
\cite{MR0202115}
は稠密部分集合上の連続関数を拡張する話．

\item 
論文
\cite{MR0036503}
はコンパクト開位相のコンパクト部分集合の話をしている気がする．何かこの論文は大事な気がする．


\item 
E. Michael
は有名な連続選択子定理の論文で，
連続選択子の理論の言葉でパラコンパクト性を特徴付けた．
論文
\cite{MR0063648}
では可算パラコンパクトと，
メタコンパクトを
選択子の理論の言葉で特徴付けしている．


\item 
論文
\cite{sierpinski1931crible}
はフランス語なのでよくわからない．
無理数の空間の話をしていて，universalな話をしているようだが，
0次元可分距離空間を埋め込めるみたいな話なのか，
それとも普遍ボレル集合みたいな話なのかはよくわかんなかった．

\item 
論文
\cite{MR0025547}
はバナッハの論文で
測度論の話をしているのはわかるが細かいところは何やっているのかフランス語だからわからなかった．

\item 
論文
\cite{Kuratowski1932}
は閉集合のなす空間(つまり超空間)
の上の半連続関数の話をしているようだが，
フランス語なのでよくわからない．
\end{enumerate}




\section*{参考文献たち}

\renewcommand{\refname}{和文参考文献}

\begin{thebibliography}{99}
\bibitem[和1]{wa1}
鈴木真治「カントールによらない実数の非可算性の証明 その発案の経緯と現在への影響をめぐって」, 
第２５回数学史シンポジウム, 
2015.


\bibitem[和２]{wa2}
松下慎也  
「On the relation between the Takahashi fixed point theorem and the Fan-Knaster-Kuratowski-Mazurkiewicz theorem」, 
数理解析研究所講究録, 
1643巻 (2009)
203--206. 

\bibitem[和３]{wa3}
玉野研一
「A space with a countable network and different dimensions --Delistathis and Watoson's example」, 
数理解析研究所講究録, 
1643巻 (1999)
61--68. 

\bibitem[和4]{wa4}
玉野研一
「積の正規性とelementary submodel」, 
数理解析研究所講究録, 
832巻 (1993)
73--79. 


\end{thebibliography}



\renewcommand{\refname}{一般の参考文献}





%READ!
%myplain is the same to amsplain style. 
\nocite{*}
\bibliographystyle{myplaindoidoi}
\bibliography{sonota.bib}



\end{document}